\begin{pastexam}[label={px:nsysu-113-dmalg-07}]{中山資工所資安組 113 學年度 第 7 題}
（設計、分析）

(Algorithm points) {[10\%]} Please express the algorithm of Quick Sort and analyze its time complexity in
detail.
\end{pastexam}

\begin{pastsoln}
\textbf{演算法}（資結書第~14.5~節）。 快速排序是分治法：選一個基準 $x$，把陣列分割成「$\le x$」「$x$」「$> x$」三段，
基準就落在它最後的位置；再遞迴排左右兩段，不需要合併。 用 Hoare 分割（取第一個數當基準）：
\begin{pseudo}{快速排序（Hoare 分割）}
\Procedure{Quicksort}{$A, p, r$}
    \If{$p < r$}
        \State $q \gets$ \Call{Hoare-Partition}{$A, p, r$}
        \State \Call{Quicksort}{$A, p, q$}
        \State \Call{Quicksort}{$A, q + 1, r$}
    \EndIf
\EndProcedure
\Function{Hoare-Partition}{$A, p, r$}
    \State $x \gets A[p]$ \Comment{基準}
    \State $i \gets p$；$j \gets r$
    \While{\TRUE}
        \While{$A[i] < x$}
            \State $i \gets i + 1$
        \EndWhile
        \While{$A[j] > x$}
            \State $j \gets j - 1$
        \EndWhile
        \If{$i \ge j$}
            \State \Return $j$
        \EndIf
        \State 交換 $A[i]$ 和 $A[j]$
        \State $i \gets i + 1$；$j \gets j - 1$
    \EndWhile
\EndFunction
\end{pseudo}
\proc{Hoare-Partition} 回傳 $j$ 時，$A[p..j]$ 的每個數都 $\le$ $A[j+1..r]$ 的每個數，而且 $p \le j < r$，兩邊都比原來
短，遞迴一定會結束（資結書虛擬碼~14.10、14.11）。 注意這個版本遞迴的是 $A[p..q]$，不是 $A[p..q-1]$：Hoare 分割不保證
基準落在 $q$。 例：$\langle 5, 3, 8, 1, 9, 2\rangle$ 以 $5$ 為基準，$i$、$j$ 從兩端往中間走，先交換 $5 \leftrightarrow 2$、
再交換 $8 \leftrightarrow 1$，得 $\langle 2, 3, 1, 8, 9, 5\rangle$，回傳 $j = 3$，左段 $\langle 2, 3, 1\rangle$ 都 $\le$ 右段
$\langle 8, 9, 5\rangle$。

\textbf{時間分析}。 一次分割掃過整段，$n$ 個數花 $\Theta(n)$。 切成 $q$ 和 $n - q$ 個時
$T(n) = T(q) + T(n - q) + \Theta(n)$。
\begin{itemize}
\item \textbf{最壞}：每次都切成 $1$ 和 $n - 1$（例如輸入已經排好，基準每次都是最小值）。
$T(n) = T(n-1) + \Theta(n) = \Theta(n) + \Theta(n-1) + \cdots = \Theta(n^2)$。 上界也成立：$q$ 怎麼選，
$q^2 + (n-q)^2 \le 1 + (n-1)^2$，用代換法可證 $T(n) \le cn^2$，所以最壞就是 $\Theta(n^2)$。
\item \textbf{最好}：每次切成兩半，$T(n) = 2T(n/2) + \Theta(n) = \Theta(n \lg n)$（主定理情況 2，資結書定理~2.4）。
只要每次至少切出固定的比例（例如 $1:9$），遞迴樹高度仍是 $O(\lg n)$、每層工作 $O(n)$，也是 $\Theta(n \lg n)$。
\item \textbf{平均}：輸入的 $n!$ 種排列等可能（或隨機選基準），基準的名次均勻分布。 比較次數 $C(n)$ 滿足
\[
C(n) = (n - 1) + \frac{2}{n}\sum_{q=0}^{n-1} C(q)
\]
（以基準落在 $q + 1$ 的 Lomuto 版寫，兩種分割的階數一樣；$C(0) = C(1) = 0$）。 兩邊乘 $n$，再減去 $n - 1$
的同一個式子：
\[
nC(n) - (n-1)C(n-1) = 2(n-1) + 2C(n-1) \;\Longrightarrow\; nC(n) = (n+1)C(n-1) + 2(n-1).
\]
除以 $n(n+1)$，並用 $\frac{2(n-1)}{n(n+1)} = \frac{4}{n+1} - \frac{2}{n}$：
\[
\frac{C(n)}{n+1} = \frac{C(n-1)}{n} + \frac{4}{n+1} - \frac{2}{n}.
\]
從 $2$ 疊加到 $n$（$C(1) = 0$），$\frac{C(n)}{n+1} = 4\bigl(H_{n+1} - \frac{3}{2}\bigr) - 2(H_n - 1) = 2H_n + \frac{4}{n+1} - 4$，
所以 $C(n) = 2(n+1)H_n - 4n \approx 1.39\, n \lg n$（$H_n$ 是調和數，$H_n \approx \ln n$；也可用資結書定理~14.3 的指示
變數法）。 平均 $\Theta(n \lg n)$。
\end{itemize}
空間：原地排序，額外空間是遞迴堆疊，平均 $\Theta(\lg n)$、最壞 $\Theta(n)$。 隨機選基準讓任何輸入的期望時間都是
$O(n \lg n)$，最壞情況不再對應到特定的輸入。
\end{pastsoln}
